\section{The homogeneous integer conclusion}
\label{im-section}

The integral map in Lemma~\ref{q13-integral-map} has only the
nonzero difference $c=a-b$ as a denominator condition. It therefore
gives the following stronger domain statement, beyond positive primitive
seeds.

\begin{corollary}[All integer points of the thirteen square class]
\label{im-all-integers}
For arbitrary integers $a,b,s$,
\[
 F(a,b)=13s^2\quad\Longleftrightarrow\quad
 a=b\ \text{ and }\ (s=a^2\ \text{ or }\ s=-a^2).
\]
\end{corollary}
\begin{proof}
If $a\ne b$, Lemma~\ref{q13-integral-map} contradicts
Lemma~\ref{q13-group}. Otherwise $F(a,a)=13a^4$, and
$(s-a^2)(s+a^2)=0$ gives the two signs. The converse follows by
substitution, including $a=b=s=0$.
\end{proof}

\section{Formal scope of the square gap and integral bridge}
\label{seventh-formal-scope}

Two new modules contain 23 declarations accepted by a fresh scoped Lean
4.32.0 build. The same build recompiles 65 unchanged dependency
declarations, with warnings treated as errors and an axiom inventory for
every declaration. Only \texttt{propext}, \texttt{Classical.choice} and
\texttt{Quot.sound} occur. This is a scoped build, not a full repository
build.

\texttt{SquareGapArithmetic} contains 16 declarations. It proves the
actual polynomial gap, transfers the modulo-eight residue table to all
primitive integer seeds, establishes positivity and the required integer
square separation, and concludes, for arbitrary signed integer $q$,
\[
 a,b>0,\quad \gcd(a,b)=1,\quad F(a,b)=q^2
 \quad\Longrightarrow\quad
 2a<3b^3,\quad 2b<3a^3,\quad a,b>1.
\]
The gap is derived from the actual power equation; it is not an extra
hypothesis. The sharper residue-dependent coefficient $2r$ in the
ordinary argument is not claimed as the conclusion of this module.

\texttt{ThirteenMapArithmetic} contains seven declarations: the three
integral identities above, the actual weighted curve equation,
nonvanishing of $U$, construction of a nontrivial integral point, and
an explicitly conditional diagonal consequence. The final theorem takes
the following antecedent as an ordinary theorem parameter:
\[
 \forall U,V,c\in\mathbb Z,\quad c\ne0,\quad
 V^2=U^3-13U^2c^2-507Uc^4\ \Longrightarrow\ U=0.
\]
It does not declare that antecedent as an axiom or prove it. The passage
to rational coordinates and the elliptic group, isogeny descent,
Mordell--Weil and torsion arguments are not formalized by these seven
declarations. Thus the complete thirteen-square-class classification
has an independently reviewed ordinary proof, while its integral bridge
has the stated narrower formal verification.

The sieved interval estimate, effective affine-slice theorems, boundary
modularity, residual support and density arguments retain their named
external inputs. None of the new modules proves the large-prime tail,
a uniform actual-point height bound, candidate elimination in all
residual classes, or standard ABC.
